\documentclass[12pt]{article}

\usepackage{amsfonts}
\usepackage{amsmath}

\begin{document}

pag 7

\hfill

\textbf{Chapter 1. The Orbits of One-Dimensional Maps.}

\hfill

\textbf{1.1 Iteration of Functions and Examples of Dynamical Systems.}

\hfill

Dynamical systems is the study of how things change over time. Examples include
the growth of populations, the change in the weather, radioactive decay, mixing of
liquids and gases such as the ocean currents, motion of the planets, the interest in a
bank account. Ideally we would like to study these with continuously varying time
(what are called flows), but in this text we will simplify matters by only considering
discrete changes in time. For example, we might model a population by measuring
it daily. Suppose that xn is the number of members of a population on day n, where
$x_0$ is the initial population. We look for a function $f :\mathbb{R} \to \mathbb{R}$, (where $\mathbb{R}$ = set of all
real numbers), for which

\[
x_1 = f(x_0), x_2 = f(x_1) \text{ and generally } x_n = f(x_{n-1}), n= 1, 2, ...
\]

This leads to the iteration of functions in the following way:

\hfill

\textbf{Definition 1.1.1} Given $x_0 \in \mathbb{R}$, the orbit of $x_0$ under f is the set

\[
O(x_0) = {x_0, f(x_0), f^2(x_0), ...},
\]

where $f^2(x_0) = f(f(x_0)), f^3(x_0) = f(f^2(x_0))$, and continuing indefinitely, so that

\[
f^n(x) = f \circ f \circ f \circ \cdots \circ f(x); (n-times composition).
\]

Set $x_n = f^n(x_0), x_1 = f(x_0), x_2 = f^2(x_0)$, so that in general

\[
x_{n+1} = f^{n+1}(x_0) = f(f^n(x_0)) = f(x_n).
\]


More generally, f may be defined on some subinterval I of $\mathbb{R}$, but in order for the
iterates of $x \in I$ under f to be defined, we need the range of f to be contained in I,
so $f: I \to I$.

\hfill

This is what we call the iteration of one-dimensional maps (as opposed to higher
dimensional maps $f: \mathbb{R}^n \to \mathbb{R}^n$, n > 1, which will be studied briefly in a later
chapter).

\hfill

\textbf{Definition 1.1.2} A (one dimensional) dynamical system is a function $f: I \to I$
where I is some subinterval of $\mathbb{R}$.

\hfill

Given such a function f, equations of the form $x_{n+1} = f(x_n)$ are examples of
difference equations. These arise in the types of examples we mentioned above. For example in biology xn may represent the number of bacteria in a culture after n hours.
There is an obvious correspondence between one-dimensional maps and these difference
equations. For example, a difference equation commonly used for calculating
square roots:

\[
x_{n+1} = \frac{1}{2}(x_n + \frac{2}{x_n}
\]
 
corresponds to the function $f(x) = \frac{1}{2} (x + \frac{2}{x_n})$. If we start by setting $x_0 = 2$ (or in
fact any real number), and then find $x_1, x_2$ etc., we get a sequence which rapidly
approaches $\sqrt{2}$ (see page 9 of Sternberg [63]). One of the issues we examine is what
exactly is happening here.

\hfill

\textbf{Examples of Dynamical Systems 1.1.3}

\hfill

\begin{enumerate}
\item \textbf{The Trigonometric Functions} Consider the iterations of the trigonometric
functions starting with $f: \mathbb{R} \to \mathbb{R}$, f(x) = sin(x). Select $x_0 \in \mathbb{R}$ at random, e.g.,
$x_0 = 2$ and set $x_{n+1} = sin(x_n)$, n = 0, 1, 2, ... Can you guess what happens to $x_n$ as
n increases? One way to investigate this type of dynamical system is to enter 2 into
our calculator, then repeatedly press the 2nd, Answer and Sin keys (you will need
to do this many times to get a good idea - it may be easier to use Mathematica, or
some similar computer algebra system). Do the same, but replace the sine function
with the cosine function. How do we explain what appears to be happening in each
case? These are questions that we aim to answer quite soon.

\item Linear Maps Probably the simplest dynamical system (and least interesting from
a chaotic dynamical point of view), for population growth arises from the iteration of
linear maps: maps of the form $f(x) = a \cdot x$. Suppose that $x_n$ = size of a population
at time n, with the property

\[
x_{n+1} = a \cdot x_n
\]

for some constant a > 0. This is an example of a linear model for the growth of the
population.

\hfill

If the initial population is $x_0 > 0$, then $x_1 = ax_0, x_2 = ax_1 = a^2x_0$, and in general
$x_n = a^nx_0$ for n = 0, 1, 2, . . .. This is the exact solution (or closed form solution)
to the difference equation $x_{n+1} = a \cdot x_n$. Clearly f(x) = ax is the corresponding
dynamical system. We can use the solution to determine the long term behavior of
the population:

$x_n$ is very well behaved since:

\begin{enumerate}
\item  if a > 1, then $x_n \to infty$ as $n \to \infty$,

\item  if 0 < a< 1 then $x_n \to 0$ as $n \to \infty$ (i.e., the population becomes extinct),
\item if a = 1, then the population remains unchanged.
\end{enumerate}

\hfill

\textbf{3. Affine maps} These are functions $f: \mathbb{R}
 \to \mathbb{R}$ of the form f(x) = ax + b ($a \neq 0$),
for constants a and b. Consider the iterates of such maps:

\[
f^2(x) = f(ax + b) = a(ax+ b) + b = a^2x + ab + b,
\]
\[
f^3(x) = a^3x + a^2b + ab + b,
\]
\[
f^4(x) = a^4x + a^3b + a^2b + ab + b,
\]

and generally

\[
f^n(x) = a^nx + a^{n-1}b + a^{n-2}b + \cdots + ab + b.
\]


Let $x_0 \in  \mathbb{R}$ and set $x_n = f^n(x_0)$, then we have

\[
\begin{matrix}
x_n & = & a^n x_0 + (a^{n-1} + a^{n-2} + \cdots + a + 1)b \\ 
& = & a^n x_0 + b \left ( \frac{a^n - 1}{a-1} \right ), \text{ if } a \neq 0 \\
\end{matrix}
\]

or

\[
x_n = \left ( x_0 + \frac{b}{a-1} \right ) a^n + \frac{b}{1-a}, \text{ if } a \neq 1
\]

is the closed form solution in this case (here we have used the formula for the sum of
a finite geometric series:

\[
\sum_{k=0}^{n-1} r^k = \frac{r^n -1}{r-1}
\]

when $r \neq 1$). If a = 1, the solution is $x_n = x_0 + nb$ .
We can use these equations to determine the long term behavior of xn. We see
that:

\begin{enumerate}
\item if $|a| < 1$ then $a^n \to 0$ as $n \to \infty$, so that

\[
lim_{n \to \infty} x_n = \frac{b}{1-a}
\]

\item  if $a > 1$, then $lim_{n \to \infty} x_n = \infty$ for b, $x_0 > 0$,

\item  if a = 1, then $lim_n \to \infty = \infty$ if $b > 0$.
The limit won’t exist if a-1.
\end{enumerate}

\item \textbf{The Logistic Map} $L{\mu}: \mathbb{R} \to \mathbb{R}, L_{\mu}(x) = \mux(1 - x)$ was introduced to model
a certain type of population growth (see [41]). Here $\mu$ is a real parameter which
is fixed. Note that if $0 < \mu \leq 4$ then $L_{\mu}$ is a dynamical system of the interval
[0, 1], i.e. $L{\mu}: [0, 1] \to [0, 1]$). For example, when $\mu = 4, L_4(x) = 4x(1 - x)$, with
$L_4([0, 1]) = [0, 1]$ with graph given below. If $\mu > 4$, then $L_{\mu}$ is no longer a dynamical
system of [0, 1] as $L_{\mu}([0, 1])$ is not a subset of [0, 1].

\hfill

\textbf{Recurrence Relations 1.1.4}

\hfill

Many sequences can be defined recursively by specifying the first term (or two),
and then stating a general rule which specifies how to obtain the nth term from the
(n − 1)th term (or other additional terms), and using mathematical induction to see
that the sequence is “well defined” for every $n \in  \mathbb{Z}^{+}$ = {1, 2, 3, ...}. For example,
n! = n-factorial can be defined in this way by specifying 0! = 1, and n! = n · (n −1)!,
for $n \in \mathbb{Z}^{+}$. The Fibonacci sequence $F_n$, can be defined by setting

\[
F_0 = 1, F_1 = 1 and F_{n+1} = F_n + F_{n−1}, for n \in \mathbb{Z}^{+}
\]

so that $F_2 = 2, F_3 = 5$ etc.

\hfill

The orbit of a point $x_0 \in \mathbb{R}$ under a function f is then defined recursively as follows:

we are given the starting value x0, and we set

\[
x_n = f(x_{n−1}), for n \in \mathbb{Z}^{+}.
\]

The principle of mathematical induction then tells that xn is defined for every n =  0,
since it is defined for n = 0, and assuming it has been defined for k = n − 1 then
$x_n = f(x_{n−1})$ defines it for k = n.

\hfill

Ideally, given a recursively defined sequence xn, we would like to have a specific
formula for xn in terms of elementary functions (so called closed form solution), but
this is often very difficult, or impossible to achieve. In the case of affine maps and
certain logistic maps, there is a closed form solution. One can use the former to study
problems of the following type:

\hfill

\textbf{Example 1.1.5} An amount \$ T is deposited in your bank account at the end of each
month. The interest is r% per period. Find the amount A(n) accumulated at the
end of n months (Assume A(0) = T).

\hfill

\textbf{Answer}. A(n) satisfies the difference equation

\[
A(n + 1) = A(n) + A(n)r + T, where A(0) = T,
\]

or

\[
A(n + 1) = A(n)(1 + r) + T,
\]

so we set $x_0 = T, a = 1+r$ and b = T in the formula of Example 1.1.3, then the
solution is

\[
A(n) = (1+r)^nT + T(\frac{(1 + r)n − 1}{1+r-1})
\]
\[
= (1+r)^nT + \frac{T}{r}((1+r)^n -1 )
\]

pag11

\hfill


\textbf{Remark 1.1.6} It is conjectured that closed form solutions for the difference equation
arising from the logistic map are only possible when $\mu = ± 2,4$ (see Exercises 1.1 #
3 for the cases where $\mu = 2, 4$, and # 7 for the case where $\mu = -2$ and also [67] for a
discussion of this conjecture).

\hfill

\textbf{Exercises 1.1}

\hfill

\begin{enumerate}
\item If $L_{\mu}(x) = \mux(1 − x)$ is the logistic map, calculate $L_{\mu}^2$ and $L^3_{\mu}(x)$.

\item Use Example 1.1.3 for affine maps to find the solutions to the difference equations:

\begin{enumerate}
\item $x_{n+1} - \frac{x_n}{3} = 2, x_0=2
\] 

\item $x_{n+1} + 3x_n = 4, x_0 = -1.$
\end{enumerate}

\item A logistic difference equation is one of the form $x_{n+1} = \mu x_n(1-x_n)$ for some fixed
$\mu \in \mathbb{R}$. Find exact solutions to the logistic equations:

\begin{enumerate}
\item $x_{n+1} = 2x_n(1 − x_n)$. Hint: use the substitution $x_n = (1 − y_n)/2$ to transform the
equation into a simpler equation that is easily solved.
\item $x_{n+1} = 4x_n(1 - x_n)$. Hint: set $x_n = sin^2(\theta_n)$ and simplify to get an equation that
is easily solved.
\end{enumerate}

\item You borrow \$ P at r % per annum, and pay off $M at the end of each subsequent
month. Write down a difference equation for the amount owing A(n) at the end of
each month (so A(0) = P). Solve the equation to find a closed form for A(n). If
P = 100, 000, M = 1000 and r = 4, after how long will the loan be paid off?

\item If 

\[
T_{\mu}(x) =  \left\{\begin{matrix}
\mu x & \text{if} & 0 \leq x < 1/2  \\
\mu(1-x) & \text{if} & 1/2 \leq x <1  \\
\end{matrix}\right.
\]

show that $T_{\mu}$ is a dynamical system of
[0, 1] for $\mu 2 (0, 2]$

\item Let $f: \mathbb{R}  \to \mathbb{R} $ be the function defined below. For each of the intervals given,
determine whether f can be considered as a dynamical system $f : I \to  I$:

\begin{enumerate}
\item $f(x) = x^3 − 3x$,

\hfill

(i) I = [−1, 1], (ii) I = [−2, 2].

\item $f(x) = 2x^3 − 6x,$

\hfill

(i) I = [−1, 1], (ii) $I = [- \sqrt{\frac{7}{2}}, \sqrt{\frac{7}{2}}],$
 (iii) I = [−4, 4].
 
\item For the following functions, find $f^2(x), f^3(x)$ and a general formula for $f^n(x)$:

\[
(i) f(x)=x^2, (ii) f(x) = |x+1|, (iii) f(x) = \left\{\begin{matrix}
2x & \text{ if }  & 0 \leq x < 1/2 \\
2x-1 & \text{ if }  & 1/2 \leq x < 1 \\
\end{matrix}\right.
\]

\item Use mathematical induction to show that if $f(x) = \frac{2}{x+1}$, then

\[
f^n(x) = \frac{2^n(x + 2) + (-1)^n(2x − 2)}{2^n(x + 2) - (-1)^n(x - 1)}
\]

\item Show that a closed form solution to the logistic difference equation when $\mu = -2$
is given by
\end{enumerate}




\end{document}

